% Generated from locale/tr/content/first-order-logic/syntax-and-semantics/first-order-languages.tex; do not hand-edit.


\olfileid[tr]{fol}{syn}{fol}

\olsection{Birinci Dereceden Diller}


Birinci dereceden mantığın ifadeleri; \emph{!!{variable}s},
\emph{!!{constant}s}, \emph{!!{predicate}s} ve kimi zaman
\emph{!!{function}s} içeren temel bir söz varlığından kurulur. Bunlardan,
mantıksal bağlaçlar, niceleyiciler ve parantezlerle virgüller gibi
noktalama simgeleriyle birlikte, \emph{terimler} ve
\emph{!!{formula}s} oluşturulur.

\begin{explain}
Gayriresmî olarak, !!{predicate}s özelliklerin ve bağıntıların,
!!{constant}s tek tek nesnelerin, !!{function}s ise gönderimlerin
adlarıdır. !!{identity}~$\eq$ dışında bunlar \emph{mantıksal olmayan
simgeler}dir ve birlikte bir dil oluştururlar. Her birinci dereceden
dil~$\Lang L$, mantıksal olmayan simgeleriyle belirlenir. En genel
durumda $\Lang L$, her türden sonsuz sayıda simge içerir.
\end{explain}

Genel durumda birinci dereceden mantıkta aşağıdaki simgeleri kullanırız:

\begin{enumerate}
\item Mantıksal simgeler
\begin{enumerate}
\item Mantıksal bağlaçlar:
  \startycommalist
  \iftag{prvNot}{\ycomma $\lnot$ (değilleme)}{}%
  \iftag{prvAnd}{\ycomma $\land$ (ve bağlacı)}{}%
  \iftag{prvOr}{\ycomma $\lor$ (veya bağlacı)}{}%
  \iftag{prvIf}{\ycomma $\lif$ (!!{conditional})}{}%
  \iftag{prvIff}{\ycomma $\liff$ (!!{biconditional})}{}%
  \iftag{prvAll}{\ycomma $\lforall$ (evrensel niceleyici)}{}%
  \iftag{prvEx}{\ycomma $\lexists$ (varlıksal niceleyici)}{}.
\tagitem{prvFalse}{!!{falsity} için önerme sabiti~$\lfalse$.}{}
\tagitem{prvTrue}{!!{truth} için önerme sabiti~$\ltrue$.}{}
\item İki yerli !!{identity}~$\eq$.
\item Bir !!{denumerable}s küme oluşturan !!{variable}s: $\Obj v_0$, $\Obj v_1$, $\Obj
  v_2$, \dots
\end{enumerate}
\item Birinci dereceden mantığın \emph{standart dilini} oluşturan
  mantıksal olmayan simgeler
\begin{enumerate}
\item Her $n>0$ için bir !!{denumerable}s küme oluşturan $n$-yerli !!{predicate}s: $\Obj
  A^n_0$, $\Obj A^n_1$, $\Obj A^n_2$, \dots
\item Bir !!{denumerable}s küme oluşturan !!{constant}s: $\Obj c_0$, $\Obj c_1$, $\Obj
  c_2$, \dots.
\item Her $n>0$ için bir !!{denumerable}s küme oluşturan $n$-yerli !!{function}s:
  $\Obj f^n_0$, $\Obj f^n_1$, $\Obj f^n_2$, \dots
\end{enumerate}
\item Noktalama işaretleri: (, ) ve virgül.
\end{enumerate}

Tanımlarımızın ve sonuçlarımızın çoğu birinci dereceden mantığın tam
standart dili için ifade edilecektir. Bununla birlikte, uygulamaya bağlı
olarak dili yalnızca belirli !!{predicate}s, !!{constant}s ve
!!{function}s içerecek biçimde de kısıtlayabiliriz.

\begin{ex}
Aritmetiğin dili~$\Lang L_A$ tek bir iki yerli !!{predicate}~$<$, tek
bir !!{constant}~$\Obj 0$, bir tane bir yerli !!{function}~$\prime$ ve
sayıları iki olan, iki yerli !!{function}{}u içerir:~$+$ ile~$\times$.
\end{ex}

\begin{ex}
Küme teorisinin dili~$\Lang L_Z$ yalnızca tek bir iki yerli
!!{predicate}~$\in$ içerir.
\end{ex}

\begin{ex}
Sıralamaların dili~$\Lang L_\le$ yalnızca iki yerli
!!{predicate}~$\le$ içerir.
\end{ex}

Bunlar da birer uzlaşımdır: resmî olarak yalnızca takma adlardır;
örneğin $<$, $\in$ ve $\le$, $\Obj A^2_0$ için, $\Obj 0$, $\Obj c_0$ için,
$\prime$, $\Obj f^1_0$ için, $+$, $\Obj f^2_0$ için ve $\times$,
$\Obj f^2_1$ için takma addır.

\iftag{defNot,defOr,defAnd,defIf,defIff,defTrue,defFalse,defEx,defAll}{%
Yukarıda tanıtılan ilkel bağlaçlara ve
\iftag{notprvEx,notprvAll}{niceleyiciye}{niceleyicilere} ek olarak
aşağıdaki \emph{tanımlı} simgeleri de kullanırız:
\startycommalist
  \iftag{defNot}{\ycomma $\lnot$ (değilleme)}{}%
  \iftag{defAnd}{\ycomma $\land$ (ve bağlacı)}{}%
  \iftag{defOr}{\ycomma $\lor$ (veya bağlacı)}{}%
  \iftag{defIf}{\ycomma $\lif$ (!!{conditional})}{}%
  \iftag{defIff}{\ycomma $\liff$ (!!{biconditional})}{}%
  \iftag{defAll}{\ycomma $\lforall$ (evrensel niceleyici)}{}%
  \iftag{defEx}{\ycomma $\lexists$ (varlıksal niceleyici)}{}%
  \iftag{defFalse}{\ycomma !!{falsity}~$\lfalse$}{}%
  \iftag{defTrue}{\ycomma !!{truth}~$\ltrue$}}{}.

\begin{tagblock}{defNot,defOr,defAnd,defIf,defIff,defTrue,defFalse,defEx,defAll}
\begin{explain}
Tanımlı bir simge resmen dilin bir parçası değildir; gayriresmî bir
kısaltma olarak sunulur. Böylece yalnızca ilkel simgeleri kullansaydık
oldukça uzayacak formülleri kısaltabiliriz. Bu açıkça bir üstünlüktür.
Fakat daha büyük üstünlük, ispatların kısalmasıdır. Bir simge ilkelse
ispatlarda ayrı olarak ele alınmalıdır. Dolayısıyla ilkel simge sayısı
arttıkça ispatlarımız da uzar.
\end{explain}
\end{tagblock}

% Alternate symbols

\begin{intro}
Yukarıda kullandığımız terimlerden ve simgelerden farklı olanlarına
aşina olabilirsiniz. Mantık metinlerinde (ve derslerde) ``değilleme''
için yaygın olarak $\sim$, $\neg$ veya~!\ kullanılır; ``ve bağlacı'' içinse
$\wedge$, $\cdot$ veya $\&$ kullanılır. ``Koşul'' ya da ``gerektirme''
için yaygın simgeler $\rightarrow$, $\Rightarrow$ ve $\supset$'dir.
\iftag{prvIff,defIff}{``İki yönlü koşul'', ``iki yönlü gerektirme''
  ya da ``(maddi) eşdeğerlik'' için kullanılan simgeler
  $\leftrightarrow$, $\Leftrightarrow$ ve $\equiv$'dir.}{}
\iftag{prvFalse,defFalse}{$\lfalse$ simgesine çeşitli kaynaklarda
  ``yanlışlık'', ``falsum'', ``saçmalık'' ya da ``alt'' denir.}{}
\iftag{prvTrue,defTrue}{$\ltrue$ simgesine çeşitli kaynaklarda
  ``doğruluk'', ``verum'' ya da ``üst'' denir.}{}

Latin alfabesinin başındaki küçük harfleri (ör. $a$, $b$, $c$)
!!{constant}s için (bunlara bazen ad da denir), sonundaki küçük harfleri
(ör. $x$, $y$, $z$) ise !!{variable}s için kullanmak gelenektir.
Niceleyiciler !!{variable}s{}le birleşir; örneğin $x$ ile. Evrensel
niceleyicinin gösterim çeşitleri $\forall x$, $(\forall x)$, $(x)$, $\Pi x$,
$\bigwedge_x$; varlıksal niceleyicininkiler ise $\exists x$, $(\exists
x)$, $(Ex)$, $\Sigma x$, $\bigvee_x$ biçimlerindedir.
\end{intro}

\begin{explain}
Bütün önerme işleçlerini ve her iki niceleyiciyi dilin ilkel simgeleri
olarak ele alabiliriz. Bunun yerine daha küçük bir ilkel simge dağarcığı
seçip öteki !!{operator}s için tanımlı simge muamelesi yapabiliriz. Doğruluk
fonksiyonları bakımından tam Boole işleci kümeleri arasında
$\{ \lnot, \lor \}$, $\{ \lnot, \land \}$ ve $\{ \lnot,
\lif\}$ bulunur---bunların her biri, anlatım bakımından tam bir birinci
dereceden dil elde etmek için niceleyicilerden herhangi biriyle
birleştirilebilir.

Başka !!{operator}s arasından ikisini de biliyor olabilirsiniz: Sheffer çizgisi~$|$
(adını Henry Sheffer'dan alır) ve Quine hançeri olarak da bilinen
Peirce oku~$\downarrow$. Bu işleçler sırasıyla olağan ``ve-değil'' ve
``veya-değil'' okumalarıyla tek başlarına doğruluk fonksiyonları
bakımından tamdır.
\end{explain}

